\documentclass[12pt,a4paper]{article}

\usepackage{ctex}
\usepackage{geometry}
\usepackage{xcolor}
\usepackage{listings}
\usepackage{graphicx}
\usepackage{float}
\usepackage{hyperref}
\usepackage{enumitem}

\geometry{margin=2.5cm}
\hypersetup{
  colorlinks=true,
  linkcolor=blue,
  urlcolor=blue
}

\lstdefinestyle{yamlstyle}{
  basicstyle=\ttfamily\small,
  breaklines=true,
  frame=single,
  columns=fullflexible,
  keepspaces=true,
  backgroundcolor=\color{gray!8}
}

\title{CI/CD 作業報告}
\author{學號：314551148 \\ 姓名：請填入姓名}
\date{\today}

\begin{document}

\maketitle

\section{CI Pipeline 說明}

本作業以 \texttt{cicd-lab} 專案為基礎，新增
\texttt{.github/workflows/ci\_314551148.yaml}，並設定在任何 branch push 時自動執行 CI pipeline。

Pipeline 主要包含以下檢查：

\begin{itemize}[noitemsep]
  \item TypeScript typecheck
  \item Prettier format check
  \item Vitest test
  \item 測試結果摘要與 JUnit XML artifact
\end{itemize}

\subsection{Workflow 主要內容}

\begin{lstlisting}[style=yamlstyle,language={}]
name: CI 314551148

on:
  push:
    branches:
      - '**'

permissions:
  contents: read

concurrency:
  group: ci-314551148-${{ github.ref }}
  cancel-in-progress: true

jobs:
  checks:
    name: Typecheck, format, and test
    runs-on: ubuntu-latest

    steps:
      - name: Checkout repository
        uses: actions/checkout@v5

      - name: Setup Node.js
        uses: actions/setup-node@v5
        with:
          node-version: '22'
          cache: npm

      - name: Install dependencies
        run: npm ci

      - name: TypeScript typecheck
        run: npm run typecheck

      - name: Prettier check
        run: npm run format:check

      - name: Run tests with JUnit report
        run: |
          mkdir -p reports
          npm test -- --reporter=default --reporter=junit --outputFile.junit=reports/vitest-junit.xml

      - name: Publish test result summary
        if: always()
        run: |
          node <<'NODE'
          const fs = require('node:fs');
          const path = 'reports/vitest-junit.xml';
          const summaryPath = process.env.GITHUB_STEP_SUMMARY;

          if (!summaryPath) {
            process.exit(0);
          }

          if (!fs.existsSync(path)) {
            fs.appendFileSync(
              summaryPath,
              [
                '## Test Results',
                '',
                'JUnit report was not generated. Check the test step logs for the failure reason.',
                ''
              ].join('\n')
            );
            process.exit(0);
          }

          const xml = fs.readFileSync(path, 'utf8');
          const attrs = Object.fromEntries(
            [...(xml.match(/<testsuites\b[^>]*>/)?.[0] ?? '').matchAll(/(\w+)="([^"]*)"/g)].map(
              ([, key, value]) => [key, value]
            )
          );

          const rows = [
            ['Total', attrs.tests ?? '0'],
            ['Failures', attrs.failures ?? '0'],
            ['Errors', attrs.errors ?? '0'],
            ['Duration (seconds)', attrs.time ?? '0']
          ];

          const markdown = [
            '## Test Results',
            '',
            '| Metric | Value |',
            '| --- | ---: |',
            ...rows.map(([label, value]) => `| ${label} | ${value} |`),
            '',
            'JUnit report path: `reports/vitest-junit.xml`',
            ''
          ].join('\n');

          fs.appendFileSync(summaryPath, markdown);
          NODE

      - name: Upload test report artifact
        if: always()
        uses: actions/upload-artifact@v7
        with:
          name: vitest-junit-report
          path: reports/vitest-junit.xml
          if-no-files-found: warn
\end{lstlisting}

\subsection{設計說明}

\begin{itemize}
  \item \texttt{npm ci}：依照 \texttt{package-lock.json} 安裝固定版本套件，確保本機與 GitHub Actions 的相依套件一致。
  \item \texttt{npm run typecheck}：執行 TypeScript 型別檢查，不產生編譯輸出。
  \item \texttt{npm run format:check}：使用 Prettier 檢查專案格式，確保程式碼風格一致。
  \item \texttt{npm test}：使用 Vitest 執行測試，並額外輸出 \texttt{reports/vitest-junit.xml}。
  \item \texttt{Publish test result summary}：讀取 JUnit XML，將測試總數、失敗數、錯誤數與執行時間寫入 GitHub Actions job summary。
  \item \texttt{Upload test report artifact}：將 JUnit XML 上傳成 artifact，方便在 GitHub Actions 結果頁下載。
\end{itemize}

GitHub Actions 中任一步驟只要回傳非 0 exit code，整個 job 就會被標記為 failed。因此 TypeScript、Prettier 或測試任一檢查失敗時，CI pipeline 會正確顯示失敗。

\section{CI 執行結果截圖}

Push 到 GitHub 後，進入 repository 的 Actions 頁面，開啟 \texttt{CI 314551148} workflow，確認 pipeline 成功執行。

\begin{figure}[H]
  \centering
  \fbox{\parbox[c][7cm][c]{0.9\textwidth}{\centering 請在此放入 GitHub Actions 成功執行截圖}}
  \caption{CI pipeline 成功執行結果}
\end{figure}

GitHub Actions 結果頁連結：\url{請填入連結}

\section{失敗案例說明}

為了驗證 CI 會在錯誤發生時失敗，我故意製造測試失敗案例。修改 \texttt{test/app.test.ts} 中 \texttt{/health} endpoint 的預期結果，例如將：

\begin{lstlisting}[style=yamlstyle,language={}]
expect(response.json()).toEqual({ status: 'ok' });
\end{lstlisting}

改成：

\begin{lstlisting}[style=yamlstyle,language={}]
expect(response.json()).toEqual({ status: 'fail' });
\end{lstlisting}

此時 API 實際回傳仍然是 \texttt{\{ status: 'ok' \}}，但測試預期值被改成 \texttt{\{ status: 'fail' \}}，因此 Vitest 會失敗，GitHub Actions 的 \texttt{Run tests with JUnit report} 步驟也會顯示 failed。

\begin{figure}[H]
  \centering
  \fbox{\parbox[c][7cm][c]{0.9\textwidth}{\centering 請在此放入 GitHub Actions 失敗執行截圖}}
  \caption{CI pipeline 失敗案例}
\end{figure}

修正方式是將測試預期值改回：

\begin{lstlisting}[style=yamlstyle,language={}]
expect(response.json()).toEqual({ status: 'ok' });
\end{lstlisting}

修正後重新 push，GitHub Actions pipeline 會再次成功通過。

\section{結論}

本次作業完成了基本 CI pipeline 設計，包含 TypeScript 型別檢查、Prettier 格式檢查、Vitest 測試、測試結果摘要與 artifact 上傳。此設計能在 push 時自動檢查專案品質，並且在任一檢查失敗時明確讓 pipeline 失敗，符合 CI 的基本要求。

\end{document}
